\section{Hash polinomial}

\begin{frame}{Hash polinomial}
  Cada caracter es un entero entre \(1\) y \(|\Sigma|\).
  Si el \(0\) fuera un caracter, las strings \(0\), \(00\), \(000\) tendrían todas hash \(0\).
  Elegimos un entero \(b > |\Sigma|\) y un primo \(p\) grande.
  \[
    h(s) = \sum_{i=0}^{n-1} s[i]\, b^{i} \pmod{p},
    \qquad n = |s|.
  \]

  \medskip
  Si \(h(s_1) = h(s_2)\), tratamos a \(s_1\) y \(s_2\) como iguales.

  \medskip
  Cuando \(s_1 = s_2\), los hashes coinciden.
  Cuando \(s_1 \neq s_2\), coinciden muy pocas veces.
\end{frame}
